
\par\noindent\textit{命题~\ref{prop:7-2-312} 的证明。}
证明与命题~\ref{prop:7-2-190} 相似。
按照定义~\ref{def:7-2-63}
中 $(n',K')=(\lVert\beta\rVert,K')$
上的序关系 $<$，
归纳构造映射 $\mathfrak H_{K',\beta}$。
归纳的第一步是构造
$\mathfrak H_{1,\beta_0}$；
它与主要归纳步骤相似且更容易，
故略去。

假定对满足
$(\lVert\beta\rVert,K)<(n',K')$
的 $\beta,K$，
已经定义 $\mathfrak H_{K,\beta}$。
现对 $\lVert\beta\rVert=n'$
定义 $\mathfrak H_{K',\beta}$。
式 (7.2.314.3) 作为滤过
$A_{n',K'}$ 同态的等式成立，
等价于
\begin{equation}
\begin{aligned}
\mathfrak H_{K',\beta}\circ
\mathfrak h_{1,\beta_0}^{\otimes K'}
&=(\mathfrak H_0)_{K',\beta}\\
&\quad-
\sum \mathfrak H_{\ell,\beta'}\circ
(\mathfrak h_{k_1,\beta_1}\otimes\cdots\otimes
\mathfrak h_{k_\ell,\beta_\ell})
\circ\Delta^{\ell-1}.
\end{aligned}
\tag{7.2.315}\label{eq:7-2-315}
\end{equation}
右端求和遍历
$\ell,\beta',(k_1,\beta_1),\ldots,
(k_\ell,\beta_\ell)$，
条件为
$\sum_i k_i=K'$、
$\beta'+\sum_i\beta_i=\beta$，
且 $(\ell,\beta')\ne(K',\beta)$。
因此，(7.2.315) 在
$\mathfrak h_{1,\beta_0}^{\otimes K'}$
的像上唯一决定
$\mathfrak H_{K',\beta}$。

\begin{lemma}{7.2.316}\label{lem:7-2-316}
存在映射
\[
\mathfrak H^I_{K',\beta}\colon
B_{K'}\overline C'[1]
\longrightarrow\overline{\widehat{\mathfrak C}}
\]
满足下列条件：

\par\noindent\textup{(7.2.317.1)}
在 $\mathfrak h_{1,\beta_0}^{\otimes K'}$
的像上，
$\mathfrak H^I_{K',\beta}
=\mathfrak H_{K',\beta}$，
并满足 (7.2.315)。

\begin{equation}
\overline{\operatorname{Eval}}_{s=s_0}
\circ\mathfrak H^I_{K',\beta}
=\mathfrak H^{s=s_0}_{K',\beta}.
\tag{7.2.317.2}\label{eq:7-2-317-2}
\end{equation}
\begin{equation}
\overline{\operatorname{Eval}}_{r=r_0}
\circ\mathfrak H^I_{K',\beta}
=\mathfrak H^{r=r_0}_{K',\beta}.
\tag{7.2.317.3}\label{eq:7-2-317-3}
\end{equation}
\end{lemma}

\par\noindent\textit{引理~\ref{lem:7-2-316} 的证明。}
定义
\begin{multline}
\mathfrak K=\ker\Biggl(
\bigoplus_{s_0=0,1}
\overline{\widehat{\mathfrak C}}_{s=s_0}
\oplus
\bigoplus_{r_0=0,1}
\overline{\widehat{\mathfrak C}}_{r=r_0}\\
\xrightarrow{\overline B}
\bigoplus_{r_0=0,1}
\bigoplus_{s_0=0,1}
\overline C_{s=s_0,r=r_0}
\Biggr).
\tag{7.2.318}\label{eq:7-2-318}
\end{multline}
由 (7.2.313.1) 可知，
\[
\bigl(
\mathfrak H^{s=0}_{K',\beta},
\mathfrak H^{s=1}_{K',\beta},
\mathfrak H^{r=0}_{K',\beta},
\mathfrak H^{r=1}_{K',\beta}
\bigr)
\in\operatorname{Hom}
(B_{K'}\overline C'[1],\mathfrak K).
\]
由 $\mathfrak h$ 的分裂单射性假设，
$\mathfrak h_{1,\beta_0}^{\otimes K'}$
的像是
$B_{K'}\overline C'$ 的一个直和项。
此外，评价映射诱导满射
$\overline{\widehat{\mathfrak C}}\to\mathfrak K$。
引理~\ref{lem:7-2-316} 得证。
\hfill$\square$

由下式定义障碍类
$o_{K',\beta}(\mathfrak H)
\in\operatorname{Hom}
(B_{K'}\overline C'[1],
\overline{\widehat{\mathfrak C}})$：
\[
\begin{aligned}
o_{K',\beta}(\mathfrak H)
+\delta_1(\mathfrak H^I_{K',\beta})
&=-\sum\mathfrak m_{\ell,\beta'}\circ
(\mathfrak H_{k_1,\beta_1}\otimes\cdots\otimes
\mathfrak H_{k_\ell,\beta_\ell})
\circ\Delta^{\ell-1}\\
&\quad+
\sum\mathfrak H_{k_1,\beta_1}\circ
\widehat{\mathfrak m}_{k_2,\beta_2}.
\end{aligned}
\]
第一个求和遍历
$\ell,\beta',(k_1,\beta_1),\ldots,
(k_\ell,\beta_\ell)$，
满足
$\sum_i k_i=K'$、
$\beta'+\sum_i\beta_i=\beta$，
且 $(\ell,\beta')\ne(1,\beta_0)$。
第二个求和遍历
$k_1,k_2,\beta_1,\beta_2$，
满足
$k_1+k_2=K'+1$、
$\beta_1+\beta_2=\beta$，
且 $(k_2,\beta_2)\ne(1,\beta_0)$。
我们有如下结论。

\begin{lemma}{7.2.319}\label{lem:7-2-319}
\begin{equation}
o_{K',\beta}(\mathfrak H)\circ
\mathfrak h_{1,\beta_0}^{\otimes K'}=0.
\tag{7.2.320.1}\label{eq:7-2-320-1}
\end{equation}
\begin{equation}
o_{K',\beta}(\mathfrak H)
\in\operatorname{Hom}
(B_{K'}\overline C'[1],\mathfrak K).
\tag{7.2.320.2}\label{eq:7-2-320-2}
\end{equation}
\begin{equation}
\delta_1\bigl(o_{K',\beta}(\mathfrak H)\bigr)=0.
\tag{7.2.320.3}\label{eq:7-2-320-3}
\end{equation}
\end{lemma}

\par\noindent\textit{引理~\ref{lem:7-2-319} 的证明。}
先证明 (7.2.320.1)。
利用 (7.2.315)，以及
$\mathfrak H_0$、$\mathfrak h$
都是滤过 $A_{n',K'}$ 同态，得到
\begin{equation}
\begin{aligned}
&\sum
\mathfrak H_{\ell',\beta'}\circ
\widehat{\mathfrak m}_{\ell'',\beta''}
\circ(\mathfrak h_{k_1,\beta_1}
\otimes\cdots\otimes\mathfrak h_{k_\ell,\beta_\ell})
\circ\Delta^{\ell-1}\\
&\quad=
\sum
\mathfrak m_{\ell,\beta'}\circ
(\mathfrak H_{k_1,\beta_1}
\otimes\cdots\otimes\mathfrak H_{k_\ell,\beta_\ell})\\
&\qquad\quad\circ
\bigotimes_{j=1}^{\ell}
\bigl(\mathfrak h_{m_{j,1},\beta_{j,1}}
\otimes\cdots\otimes
\mathfrak h_{m_{j,k_j},\beta_{j,k_j}}\bigr)
\circ\Delta^{\sum_j k_j-1}.
\end{aligned}
\tag{7.2.321}\label{eq:7-2-321}
\end{equation}
左端求和遍历
$\ell,\ell',\ell'',\beta',\beta'',
k_i,\beta_i$，
满足 $\ell'+\ell''=\ell+1$、
$\sum_i k_i=K'$、
$\beta'+\beta''+\sum_i\beta_i=\beta$。
右端求和遍历\translatornote{原书式 (7.2.321)
的右端张量积末因子下标排作
$m_{k_j},\beta_{j,m_{k_j}}$；
这里按下文求和指标
$m_{j,i},\beta_{j,i}$ 写为
$m_{j,k_j},\beta_{j,k_j}$。}
$\ell,\beta',k_i,\beta_i,
m_{j,i},\beta_{j,i}$，
满足
$\sum_{i,j}m_{j,i}=K'$、
$\beta'+\sum_j\beta_j+
\sum_{j,i}\beta_{j,i}=\beta$。

与引理~\ref{lem:7-2-193} 的证明相同，
式 (7.2.321) 中凡含有至少一个
$\mathfrak h_{a,\beta(a)}$、
且 $(a,\beta(a))\ne(1,\beta_0)$
的项，
都由归纳假设相消；
该假设是 $\mathfrak H$
在 $(n,K)<(n',K')$ 时
已为滤过 $A_{n,K}$ 同态。
这就推出 (7.2.320.1)。
(7.2.320.2) 由
(7.2.317.2)、(7.2.317.3)
及归纳假设得到。
(7.2.320.3)\translatornote{原书此处误写
“(7.2.317.3)”，
但本段末项所需的是障碍类的闭性
(7.2.320.3)。}
由定义及归纳假设得到。
引理~\ref{lem:7-2-319} 得证。
\hfill$\square$

由引理~\ref{lem:7-2-319} 可得，
存在
$\operatorname{Cor}\in
\operatorname{Hom}
(B_{K'}\overline C'[1],\mathfrak K)$，
使
\begin{equation}
o_{K',\beta}(\mathfrak H)
=\delta_1(\operatorname{Cor}),
\qquad
\operatorname{Cor}\circ
\mathfrak h_{1,\beta_0}^{\otimes K'}=0.
\tag{7.2.322}\label{eq:7-2-322}
\end{equation}
置
$\mathfrak H_{K',\beta}
=\mathfrak H^I_{K',\beta}
+\operatorname{Cor}$。
直接检验可知它满足所需性质。
\mbox{命题~\ref{prop:7-2-312} 得证。\(\square\)}
